\exercisegroup

\begin{enumerate}
\def\labelenumi{\arabic{enumi})}
\tightlist
\item
  \begin{enumerate}
  \def\labelenumii{\alph{enumii})}
  \tightlist
  \item
    Show that, if a Hausdorff topological vector space E over a non-discrete valued division ring K is such that every neighbourhood of 0 contains a vector subspace that is not the single point 0, then the topology of E cannot be defined by a norm. In particular, a product of an infinite sequence \((\mathrm{E}_{n})\) of Hausdorff topological vector spaces on K, none consisting of the single point 0, has a topology that cannot be defined by a norm.
  \end{enumerate}
\end{enumerate}

\begin{enumerate}
\def\labelenumi{\alph{enumi})}
\setcounter{enumi}{1}
\tightlist
\item
  Consider the product vector space \(E = K_{s}^{N}\) ; for all \(x = (\xi_{n}) \in E\) , put
\end{enumerate}

\[
| x | = \sum_ {n = 0} ^ {\infty} 2 ^ {- n} \left| \xi_ {n} \right| / (1 + | \xi_ {n} |).
\]

Show that the topology of E is defined by the distance \(d(x, y) = |x - y|\) , that \(|\lambda x| \leqslant |x|\) if \(|\lambda| \leqslant 1\) , \(|\lambda x| \leqslant |\lambda| \cdot |x|\) if \(|\lambda| \geqslant 1\) and that, for all \(x_0 \in \mathrm{E}\) , \(|\lambda x_0|\) tends to 0 with \(|\lambda|\) .

\begin{enumerate}
\def\labelenumi{\arabic{enumi})}
\setcounter{enumi}{1}
\tightlist
\item
  Let E and F be two complete, metrisable vector spaces over a non-discrete valued division ring, and let \(T_{0}\) be the topology of F. Let T be a Hausdorff topology on F, coarser than \(T_{0}\) . Show that if the linear mapping u of E in F is continuous for the topology T on F, it is still continuous for the topology \(T_{0}\) on F (use the cor. 5 of I, p.~19).
\end{enumerate}

Deduce that if \(T_{1}\) and \(T_{2}\) are two distinct topologies on a vector space E over a non-discrete valued division ring, compatible with the vector space structure of E, and for each of them E is metrisable and complete, then there does not exist a Hausdorff topology on E coarser than \(T_{1}\) and \(T_{2}\) . Give an example of two such topologies on an infinite dimensional vector space E (note that there exist bijections of E on itself such that both the bijection and its inverse are not continuous for a normed space topology on E).

\begin{enumerate}
\def\labelenumi{\arabic{enumi})}
\setcounter{enumi}{2}
\item
  Let E and F be two Hausdorff topological vector spaces over a non-discrete valued division ring; and suppose that E is metrisable and complete. Let u be a continuous linear injection of E in F, and let G be a vector subspace of \(u(\mathrm{E})\) ; suppose that there exists on G a topology T which is finer than the topology induced by that of F, is compatible with the vector space structure of G and for which G is metrisable and complete. Show that the mapping inverse to u, restricted to G, is continuous for T (use I, p.~19, cor. 5).
\item
  Let E, F be two complete metrisable vector spaces over a non-discrete valued division ring and let u be a continuous linear mapping of E in F. Show that if there exists in F a closed complementary subspace to \(u(\mathrm{E})\) , then \(u(\mathrm{E})\) is closed in F (use I, p.~19, cor. 5).
\item
  Let E and F be two complete metrisable vector spaces over a non-discrete valued division ring and let u be a linear mapping of E in F. Let N be the set of cluster points of u in F with respect to the filter of neighbourhoods of 0 in E; show that N is a closed vector subspace of F, and that, in order that u be continuous, it is necessary and sufficient that N be the single point 0 (use I, p.~19, cor. 5). Show that N is the smallest of the closed vector subspaces M of F such that, if \(\phi\) denotes the canonical homomorphism of F on F/M, then \(\phi \circ u\) is a continuous mapping of E in F/M.
\item
  Let E be a complete metrisable vector space over a non-discrete valued division ring K.
\end{enumerate}

\begin{enumerate}
\def\labelenumi{\alph{enumi})}
\tightlist
\item
  Let p be a lower semi-continuous mapping of E in the interval \([0, +\infty)\) of \(\overline{R}\) such that \(p(\lambda x) = |\lambda| \cdot p(x)\) for \(\lambda \neq 0\) in K and \(x \in E\) , such that \(p(0) = 0\) and \(p(x + y) \leqslant p(x) + p(y)\)
\end{enumerate}

for any \(x, y\) in E. Show that if \(p\) is finite in E, then \(p\) is continuous (consider the closed set B of \(x \in \mathbf{E}\) such that \(p(x) \leqslant 1\) , and use Baire's theorem).

\begin{enumerate}
\def\labelenumi{\alph{enumi})}
\setcounter{enumi}{1}
\tightlist
\item
  Let \((p_n)\) be a sequence of mappings of \(E\) in \([0, +\infty]\) satisfying the conditions of \(a\) ). Show that if none of the \(p_n\) are finite in \(E\) then there exists a point \(x \in E\) such that \(p_n(x) = +\infty\) for all \(n\) (same method as above).
\end{enumerate}

\begin{enumerate}
\def\labelenumi{\arabic{enumi})}
\setcounter{enumi}{6}
\tightlist
\item
  Let E be a complete metrisable vector space over a non-discrete valued division ring K. We say that a vector subspace M of E is paracomplete if there exists on M a complete metrisable vector space structure for which the canonical injection of M in E is continuous. a) Let M, N be two paracomplete subspaces of E such that \(M + N\) and \(M \cap N\) are closed in E. Show that M and N are closed in E. (Taking quotients by \(M \cap N\) reduces the question to the case where \(M \cap N = \{0\}\) , and we can consider then the mapping \((x, y) \mapsto x + y\) of \(M \times N\) in E).
\end{enumerate}

\begin{enumerate}
\def\labelenumi{\alph{enumi})}
\setcounter{enumi}{1}
\tightlist
\item
  Show that if E is the union of an increasing sequence of paracomplete subspaces, \((\mathbf{M}_{j})_{j\geq0}\) , then there exists an index j such that \(M_{j}=E\) . (Use Baire's theorem (GT, IX, § 5.3, th. 1) and I, p.~17, th. 1).
\end{enumerate}

\begin{enumerate}
\def\labelenumi{\arabic{enumi})}
\setcounter{enumi}{7}
\tightlist
\item
  Let E be a Banach space over a non-discrete valued division ring K. We say that a vector subspace M of E is strongly paracomplete if there exists a norm \(\|x\|_{M}\) on M for which M is a Banach space and the canonical injection of M in E is continuous.
\end{enumerate}

\begin{enumerate}
\def\labelenumi{\alph{enumi})}
\item
  Show that if \(\mathbf{M}\) and \(\mathbf{N}\) are two strongly paracomplete subspaces of \(\mathbf{E}\) , then \(\mathbf{M} + \mathbf{N}\) and \(\mathbf{M} \cap \mathbf{N}\) are also strongly paracomplete subspaces. (On \(\mathbf{M} + \mathbf{N}\) , consider the norm \(\|x\|_{\mathbf{M} + \mathbf{N}} = \inf (\|u\|_{\mathbf{M}} + \|v\|_{\mathbf{N}})\) , where the lower bound is taken over all pairs \((u, v)\) such that \(x = u + v\) , \(u \in \mathbf{M}\) and \(v \in \mathbf{N}\) .)
\item
  Let M, N be two strongly paracomplete subspaces of E such that N and \(M + N\) are closed. Show that \(\overline{M} = M + (\overline{M} \cap \overline{N})\) and \(\overline{M} \cap \overline{N} = \overline{M} \cap N\) (use exerc. 7, a)).
\end{enumerate}

\begin{enumerate}
\def\labelenumi{\arabic{enumi})}
\setcounter{enumi}{8}
\tightlist
\item
  \begin{enumerate}
  \def\labelenumii{\alph{enumii})}
  \tightlist
  \item
    Let \(a, b\) be two points of a normed space \(\mathbf{E}\) on the field \(\mathbf{R}\) . Denote by \(\delta(\mathbf{A})\) the diameter of a bounded set \(\mathbf{A}\) in \(\mathbf{E}\) (using the norm metric on \(\mathbf{E}\) ) and define inductively the sequence \((\mathbf{B}_n)_{n \geqslant 1}\) of bounded sets in \(\mathbf{E}\) satisfying the following conditions: \(\mathbf{B}_1\) is the set of those \(x \in \mathbf{E}\) such that \(\|x - a\| = \|x - b\| = \frac{1}{2}\|a - b\|\) ; for \(n > 1\) , \(\mathbf{B}_n\) is the set of those \(x \in \mathbf{B}_{n-1}\) such that \(\|x - y\| \leqslant \frac{1}{2}\delta(\mathbf{B}_{n-1})\) for all \(y \in \mathbf{B}_{n-1}\) . Show that the intersection of the \(\mathbf{B}_n\) is just the single point \(\frac{1}{2}(a + b)\) (note that \(\delta(\mathbf{B}_n) \leqslant \frac{1}{2}\delta(\mathbf{B}_{n-1})\) ).
  \end{enumerate}
\end{enumerate}

\begin{enumerate}
\def\labelenumi{\alph{enumi})}
\setcounter{enumi}{1}
\tightlist
\item
  Deduce from \(a)\) that if \(u\) is an isometry of the real Banach space \(E\) on the real Banach space \(F\) , then \(u\) is an affine linear mapping of \(E\) on \(F\) .
\end{enumerate}

